Una tira de longitud \(L\), sección constante \(A\), conductividad \(k\), densidad \(\rho\) y calor específico \(c\), todos positivos y constantes, se modela por
\[
T_t=\frac{k}{\rho c}T_{xx},\qquad 0<x<L.
\]
No hay fuentes interiores. Las caras laterales y el extremo \(x=0\) están aislados. Por \(x=L\), un calefactor entrega a la tira toda su potencia constante \(P>0\). El flujo positivo hacia \(+x\) obedece \(q=-kT_x\).

Escriba las dos condiciones de contorno con sus signos. Obtenga la variación de la energía
\[
E(t)=\rho c A\int_0^L\bigl(T(x,t)-T_{\rm ref}\bigr)\,dx,
\]
con \(T_{\rm ref}\) constante. Explique oralmente si puede alcanzarse una temperatura estacionaria bajo estas hipótesis y por qué una potencia prescrita no equivale a mantener el extremo a una temperatura prescrita. No se pide resolver la EDP.
